\documentclass[a4paper,12pt]{article}
\usepackage[utf8]{inputenc}
\usepackage{amsmath, amssymb}
\usepackage{amsthm}
\usepackage{geometry}
\usepackage{graphicx}
\usepackage[utf8]{inputenc}
\usepackage[T1]{fontenc}
\usepackage[spanish]{babel}
\usepackage{url}
\geometry{margin=1in}
\newtheorem{definition}{Definition}
\newtheorem{proposition}{Proposition}


% Datos de la tarea
\title{\textbf{Seminario Matemáticas Aplicadas II\\Bibliografia y fuentes ADP 1}}
\author{Eliud Daniel Santillan Aparicio \\ Leonardo Vigueras \\ 415069696}
\date{\today}
%\date{} % Elimina la fecha


\begin{document}

\maketitle



\subsection*{Fuentes de Datos e Información para Analizar las Dinámicas Monopolísticas de la Compañía Británica de las Indias Orientales}

\subsubsection*{Fuentes de Datos Primarias}

\begin{itemize}
    \item \textbf{Actas de la Carta (1600–1858)}
    \begin{itemize}
        \item Carta de 1600: Estableció el monopolio sobre el comercio entre el Cabo de Buena Esperanza y el Estrecho de Magallanes\cite{ref4}.
        \item Actas de 1781/1793: Renovaron el monopolio mientras formalizaban la soberanía de la Corona\cite{ref1,ref5}.
        \item Acta de 1813: Terminó el monopolio comercial indio (retuvo el té y el opio de China)\cite{ref3,ref6}.
        \item Acta de 1833: Abolió todos los monopolios comerciales\cite{ref3,ref5}.
        \item Debates Parlamentarios: Revelan motivaciones políticas (por ejemplo, necesidades fiscales, cabildeo)\cite{ref1,ref5}.
    \end{itemize}
    \item \textbf{Registros Financieros}
    \begin{itemize}
        \item Préstamos Gubernamentales: Préstamo de £1 millón en 1744 aseguró el monopolio hasta 1780\cite{ref1}.
        \item Participación en Beneficios: Los ingresos fiscales de Bengala posteriores a 1757 financiaron la industrialización británica\cite{ref2,ref6}.
        \item Datos del Comercio del Té: Las ganancias del monopolio alimentaron las tensiones coloniales (por ejemplo, la Ley del Té de 1773)\cite{ref8}.
    \end{itemize}
    \item \textbf{Estadísticas Comerciales}
    \begin{itemize}
        \item Importaciones Textiles: La tela británica en Bengala aumentó del 25\% (1811) al 93\% (1840)\cite{ref2}.
        \item Exportaciones de Materias Primas: El algodón y el opio indios dominaron los mercados mundiales\cite{ref2,ref7}.
    \end{itemize}
\end{itemize}

\subsubsection*{Marcos Analíticos}

\begin{itemize}
    \item \textbf{Economía Política}
    \begin{itemize}
        \item Estado Fiscal-Militar: La estabilidad del monopolio posterior a 1750 estuvo ligada a la Oligarquía Whig y a la mejora de las finanzas públicas\cite{ref1,ref6}.
        \item Poder de Cabildeo: Las conexiones de los directores con el Parlamento aseguraron la renovación (por ejemplo, las Actas de 1781/1793)\cite{ref1,ref5}.
    \end{itemize}
    \item \textbf{Mecanismos de Control del Mercado}
    \begin{itemize}
        \item Aranceles: Los bienes británicos entraron en la India libres de impuestos, mientras que los textiles indios enfrentaron aranceles prohibitivos en Gran Bretaña\cite{ref2,ref7}.
        \item Comercialización Forzada: Campesinos obligados a cultivar índigo y opio para la exportación\cite{ref7}.
    \end{itemize}
    \item \textbf{Impacto Global}
    \begin{itemize}
        \item Desindustrialización: Las exportaciones textiles indias colapsaron después de 1813\cite{ref2}.
        \item Comercio del Opio: El monopolio sobre las ventas de opio chinas financió las operaciones de la Compañía\cite{ref8}.
    \end{itemize}
\end{itemize}

\subsubsection*{Fuentes Secundarias Clave}

\begin{itemize}
    \item Dincecco, J. A. (2011). \textit{Fiscal Centralization, Limited Government, and Public Debt: England, 1688-1790}. NBER Working Paper No. 16831\cite{ref1,ref6}.
    \item Chaudhuri, K. N. (1978). \textit{The Trading World of Asia and the English East India Company 1660-1760}\cite{ref4}. Cambridge University Press.
    \item Dalrymple, W. (2020). Archivos sobre registros marítimos y comercio del Golfo\cite{ref3}.
\end{itemize}

\subsubsection*{Análisis Comparativo}

\begin{itemize}
    \item \textbf{Factor: Seguridad del Monopolio}
    \begin{itemize}
        \item Dinámicas Pre-1750: Renegociaciones frecuentes debido a un estado fiscal débil\cite{ref1,ref4}.
        \item Consolidación Post-1750: Renovaciones estables mediante conexiones políticas\cite{ref1,ref6}.
    \end{itemize}
    \item \textbf{Factor: Fuentes de Ingresos}
    \begin{itemize}
        \item Dinámicas Pre-1750: Ganancias comerciales.
        \item Consolidación Post-1750: Ingresos fiscales de Bengala + ventas de opio\cite{ref6,ref7}.
    \end{itemize}
    \item \textbf{Factor: Oposición}
    \begin{itemize}
        \item Dinámicas Pre-1750: Intrusos + compañías rivales\cite{ref4,ref6}.
        \item Consolidación Post-1750: Suprimida por el poder militar\cite{ref5,ref8}.
    \end{itemize}
\end{itemize}

\subsubsection*{Preguntas Sin Resolver}

\begin{itemize}
    \item Efectos en el Bienestar: ¿Superaron las ganancias del monopolio las pérdidas de eficiencia derivadas de la restricción de la competencia?\cite{ref4}
    \item Costos a Largo Plazo: ¿Cómo influyó la desindustrialización en el declive del PIB de la India en el siglo XIX?\cite{ref2,ref7}
\end{itemize}

Utilice estas fuentes para explorar cómo las rentas de monopolio, la colusión estatal y las redes comerciales globales se cruzaron en las operaciones de la Compañía.

\begin{center}
    \(\ddagger\)
\end{center}
\textbf{Bibliografía}

\textbf{Bibliografía}

\begin{itemize}
    \item \textbf{Dincecco, J. A.} (2011). \textit{Fiscal Centralization, Limited Government, and Public Debt: England, 1688-1790}. NBER Working Paper No. 16831. \url{https://www.nber.org/system/files/chapters/c13506/revisions/c13506.rev0.pdf} \cite{ref1,ref6}
    \item \textit{Economy of India under Company rule}. Wikipedia. \url{https://en.wikipedia.org/wiki/Economy_of_India_under_Company_rule} \cite{ref2}
    \item \textit{Digitised East India Company ships’ journals and related records}. British Library Blogs. \url{https://blogs.bl.uk/asian-and-african/2020/05/digitised-east-india-company-ships-journals-and-related-records.html} \cite{ref3}
    \item \textbf{Chaudhuri, K. N.} (1978). \textit{The Trading World of Asia and the English East India Company 1660-1760}. Cambridge University Press. \cite{ref4}
    \item \textit{East India Company and Raj, 1785-1858}. UK Parliament. \url{https://www.parliament.uk/about/living-heritage/evolutionofparliament/legislativescrutiny/parliament-and-empire/parliament-and-the-american-colonies-before-1765/east-india-company-and-raj-1785-1858/} \cite{ref5}
    \item \textbf{Bogart, D.} (2016). \textit{The East India Company, Privateering, and the Fiscal-Military State in Britain, 1688-1815}. Working Paper. \cite{ref6}
    \item \textit{India Under British Rule: Economic Impact}. National Institute of Open Schooling. \url{https://nios.ac.in/media/documents/secsocscicour/english/lesson-05.pdf} \cite{ref7}
    \item \textit{British East India Company}. American Battlefield Trust. \url{https://www.battlefields.org/learn/articles/british-east-india-company} \cite{ref8}
\end{itemize}
 



\end{document}